\section{How URF Unifies Domains}
\label{sec:universal-translation-spine-v1}

The Universal Translation Spine is the registry-level form of the URF unification claim. It does not assert that every domain theorem has been solved. It records the common structural shape used to compare domain surfaces.

\[
\mathrm{domain}
\longrightarrow
\mathrm{encoder}
\longrightarrow
\mathrm{rigidity\ grammar\ fragment}
\longrightarrow
\mathrm{invariant}
\longrightarrow
\mathrm{frontier\ status}.
\]

A domain is any mathematical, computational, physical, or information-theoretic setting under study. An encoder records how that domain is translated into URF data. A rigidity grammar fragment records the shared structural language. An invariant records what is preserved, obstructed, bounded, or forced. A frontier status records what is closed, conditional, restricted, or still open.

This spine makes the unification thesis auditable: different subjects are not treated as identical, but they are compared through the same tuple type.

\begin{center}
\begin{tabular}{p{0.22\linewidth}p{0.25\linewidth}p{0.25\linewidth}p{0.18\linewidth}}
\toprule
Domain & Encoder & Invariant & Frontier status \\
\midrule
Computation & local process data & width obstruction & open or conditional \\
Physics & operational constraint data & boundary rigidity & open or conditional \\
Gravity & collapse boundary data & compactness obstruction & open or conditional \\
Entropy fibers & fiber mass data & entropy balance & open or conditional \\
Temporal relaxation & dissipation data & decay obstruction & open or conditional \\
\bottomrule
\end{tabular}
\end{center}

The spine is therefore a structural unification surface. It is useful because it gives every domain the same audit fields, while preserving the boundary between registry consistency and theorem closure.

It does not prove unrestricted Chronos-RR, unrestricted H4.1/FGL, P vs NP, or any Clay problem.
